% ============================================================================
% WORKSHEET - Section 1.1: Unit Rates with Fractions
% CCSS 7.RP.A.1
% ============================================================================
\section{Unit Rates with Fractions}
\setcounter{wsNumber}{0}

\begin{worksheetSkills}
\begin{itemize}
	\item Compute unit rates when the measurements are fractions or mixed numbers
	\item Compare two situations by converting each rate to a per-$1$ value
	\item Explain why dividing by a fraction can produce a larger number
\end{itemize}
\end{worksheetSkills}

\worksheetHeader{Fraction Rate Sprint}

\begin{speedDrill}{Find Each Unit Rate}
\timerBadge{4}

For each situation, divide the first quantity by the second quantity.

\bigskip
\renewcommand{\arraystretch}{1.9}
\begin{center}
\begin{tabular}{|C{7cm}|C{4cm}|}
\hline
\textbf{Situation} & \textbf{Unit Rate} \\
\hline
$\frac{3}{4}$ cup of juice in $\frac{1}{2}$ minute & \answerBlank[3cm] \\
\hline
$\frac{5}{6}$ mile in $\frac{1}{3}$ hour & \answerBlank[3cm] \\
\hline
$\frac{7}{8}$ gallon for $\frac{1}{4}$ acre & \answerBlank[3cm] \\
\hline
$\frac{2}{3}$ yard in $\frac{5}{6}$ minute & \answerBlank[3cm] \\
\hline
$1\frac{1}{2}$ pounds in $\frac{3}{4}$ day & \answerBlank[3cm] \\
\hline
$\frac{9}{10}$ page in $\frac{3}{5}$ minute & \answerBlank[3cm] \\
\hline
\end{tabular}
\end{center}
\end{speedDrill}
\worksheetAnswer{(1)~$\frac{3}{4} \div \frac{1}{2} = \frac{3}{4} \times 2 = \frac{3}{2} = 1\frac{1}{2}$ cups per minute. (2)~$\frac{5}{6} \div \frac{1}{3} = \frac{5}{6} \times 3 = \frac{15}{6} = 2\frac{1}{2}$ miles per hour. (3)~$\frac{7}{8} \div \frac{1}{4} = \frac{7}{8} \times 4 = \frac{28}{8} = 3\frac{1}{2}$ gallons per acre. (4)~$\frac{2}{3} \div \frac{5}{6} = \frac{2}{3} \times \frac{6}{5} = \frac{12}{15} = \frac{4}{5}$ yard per minute. (5)~$1\frac{1}{2} = \frac{3}{2}$, so $\frac{3}{2} \div \frac{3}{4} = \frac{3}{2} \times \frac{4}{3} = 2$ pounds per day. (6)~$\frac{9}{10} \div \frac{3}{5} = \frac{9}{10} \times \frac{5}{3} = \frac{45}{30} = \frac{3}{2} = 1\frac{1}{2}$ pages per minute. To find a unit rate with fractions, divide and multiply by the reciprocal.}

\worksheetHeader{Unit Rate Match-Up}

\begin{matchIt}{Connect Each Situation to Its Unit Rate}

Draw a line from each situation on the left to the matching unit rate on the right.

\matchRow{$\frac{1}{2}$ mile in $\frac{1}{4}$ hour}{$\frac{4}{5}$ page per minute}
\matchRow{$\frac{3}{5}$ page in $\frac{3}{4}$ minute}{$2$ miles per hour}
\matchRow{$\frac{5}{8}$ gallon for $\frac{1}{2}$ garden bed}{$1\frac{1}{4}$ gallons per bed}
\matchRow{$\frac{7}{6}$ cups for $\frac{2}{3}$ batch}{$1\frac{3}{4}$ cups per batch}
\matchRow{$\frac{9}{10}$ yard in $\frac{3}{5}$ minute}{$1\frac{1}{2}$ yards per minute}
\end{matchIt}
\worksheetAnswer{$\frac{1}{2}$ mile in $\frac{1}{4}$ hour matches $2$ miles per hour because $\frac{1}{2} \div \frac{1}{4} = 2$. $\frac{3}{5}$ page in $\frac{3}{4}$ minute matches $\frac{4}{5}$ page per minute because $\frac{3}{5} \div \frac{3}{4} = \frac{3}{5} \times \frac{4}{3} = \frac{4}{5}$. $\frac{5}{8}$ gallon for $\frac{1}{2}$ garden bed matches $1\frac{1}{4}$ gallons per bed because $\frac{5}{8} \div \frac{1}{2} = \frac{10}{8} = \frac{5}{4}$. $\frac{7}{6}$ cups for $\frac{2}{3}$ batch matches $1\frac{3}{4}$ cups per batch because $\frac{7}{6} \div \frac{2}{3} = \frac{7}{4} = 1\frac{3}{4}$. $\frac{9}{10}$ yard in $\frac{3}{5}$ minute matches $1\frac{1}{2}$ yards per minute because $\frac{9}{10} \div \frac{3}{5} = \frac{3}{2}$.}

\worksheetHeader{Who Found the Better Rate?}

\begin{errorDetective}{Catch the Fraction Mistakes}

Each student tried to find a unit rate. Decide whether the work is correct and fix any errors.

Mia says: ``$\frac{3}{4}$ mile in $\frac{1}{2}$ hour is $\frac{3}{8}$ mile per hour.''
\errorWork{$\frac{3}{4} \div \frac{1}{2} = \frac{3}{8}$}
Correct or wrong? \answerBlank[2cm] \quad Correct rate: \answerBlank[4cm]

\bigskip

Leo says: ``$1\frac{1}{4}$ cups in $\frac{1}{2}$ batch is $\frac{5}{4}$ cups per batch.''
\errorWork{$1\frac{1}{4} = \frac{5}{4},\; \frac{5}{4} \div \frac{1}{2} = \frac{5}{4}$}
Correct or wrong? \answerBlank[2cm] \quad Correct rate: \answerBlank[4cm]

\bigskip

Zara says: ``$\frac{2}{3}$ yard in $\frac{5}{6}$ minute is $\frac{10}{18}$ yard per minute.''
\errorWork{$\frac{2}{3} \div \frac{5}{6} = \frac{2}{3} \times \frac{5}{6} = \frac{10}{18}$}
Correct or wrong? \answerBlank[2cm] \quad Correct rate: \answerBlank[4cm]
\end{errorDetective}
\worksheetAnswer{All three are wrong. (1)~Mia should multiply by the reciprocal: $\frac{3}{4} \div \frac{1}{2} = \frac{3}{4} \times 2 = \frac{3}{2} = 1\frac{1}{2}$ miles per hour. She multiplied denominators instead of dividing by $\frac{1}{2}$. (2)~Leo converted the mixed number correctly to $\frac{5}{4}$, but then forgot that dividing by $\frac{1}{2}$ doubles the amount: $\frac{5}{4} \div \frac{1}{2} = \frac{5}{4} \times 2 = \frac{10}{4} = 2\frac{1}{2}$ cups per batch. (3)~Zara used the wrong reciprocal. The correct calculation is $\frac{2}{3} \times \frac{6}{5} = \frac{12}{15} = \frac{4}{5}$ yard per minute.}

\worksheetHeader{Snack Lab Rates}

\begin{mathScenario}{Compare the Trail Mix Recipes}

Three student groups mixed snack bags for the school hike.

\bigskip
\scenarioItem{Group A}{$\frac{3}{4}$ cup raisins in $\frac{1}{2}$ batch}
\scenarioItem{Group B}{$\frac{5}{6}$ cup raisins in $\frac{2}{3}$ batch}
\scenarioItem{Group C}{$1\frac{1}{8}$ cups raisins in $\frac{3}{4}$ batch}

\bigskip

Which group used the greatest amount of raisins per batch? \answerBlank[5cm]

\bigskip

How much more raisins per batch did Group C use than Group B? \answerBlank[5cm]

\bigskip

If Group A made $3$ full batches, how many cups of raisins would they need? \answerBlank[5cm]
\end{mathScenario}
\worksheetAnswer{(1)~Group A: $\frac{3}{4} \div \frac{1}{2} = \frac{3}{2} = 1\frac{1}{2}$ cups per batch. Group B: $\frac{5}{6} \div \frac{2}{3} = \frac{5}{4} = 1\frac{1}{4}$ cups per batch. Group C: $1\frac{1}{8} = \frac{9}{8}$, and $\frac{9}{8} \div \frac{3}{4} = \frac{9}{8} \times \frac{4}{3} = \frac{36}{24} = \frac{3}{2} = 1\frac{1}{2}$ cups per batch. Groups A and C tie for the greatest unit rate at $1\frac{1}{2}$ cups per batch. (2)~Group C used $1\frac{1}{2} - 1\frac{1}{4} = \frac{1}{4}$ cup more per batch than Group B. (3)~Group A needs $1\frac{1}{2} \times 3 = 4\frac{1}{2}$ cups of raisins for $3$ batches.}

\worksheetHeader{Explain the Fraction Flip}

\begin{explainIt}{Why Can a Unit Rate Be Bigger Than Both Fractions?}

Jaden says, ``A unit rate should always be smaller than the numbers you start with, so $\frac{3}{4} \div \frac{1}{2}$ cannot equal $1\frac{1}{2}$.''
Explain why Jaden is wrong. Use the idea of ``how many halves fit into three fourths'' in your explanation.

\writeLines{6}

\bigskip

Nora compares $\frac{2}{3}$ mile in $\frac{1}{4}$ hour and $\frac{5}{6}$ mile in $\frac{1}{3}$ hour. Which unit rate is greater, and how do you know without using a calculator?

\writeLines{6}
\end{explainIt}
\worksheetAnswer{(1)~Jaden is wrong because dividing by a fraction smaller than $1$ asks how many small pieces fit into the amount. Since a half is smaller than one whole, more than one half can fit into $\frac{3}{4}$. In fact, $\frac{3}{4} \div \frac{1}{2} = \frac{3}{4} \times 2 = \frac{3}{2} = 1\frac{1}{2}$, so one and a half halves fit into three fourths. That is why the unit rate can be larger than both original fractions. (2)~First rate: $\frac{2}{3} \div \frac{1}{4} = \frac{8}{3} = 2\frac{2}{3}$ miles per hour. Second rate: $\frac{5}{6} \div \frac{1}{3} = \frac{5}{2} = 2\frac{1}{2}$ miles per hour. Since $2\frac{2}{3} > 2\frac{1}{2}$, the first unit rate is greater. Without a calculator, compare the mixed numbers or convert both to improper fractions with a common denominator.}

\worksheetDone
